% Notes for computational comparison with Euclid pipeline
% Based on:
%   - Euclid Collaboration 2025 (2501.16555) - 2PCF estimation methodology
%   - Keihänen et al. 2022 (2205.11852) - Linear Construction covariance speedup
%   - Kerscher et al. 2022 (2210.xxxxx) - Improved accuracy estimators

% KEY FINDINGS FROM EUCLID PIPELINE:

% 1. SCALE OF THE PROBLEM
%    - 20-30 million Halpha galaxies, z = 0.9-1.8
%    - Random catalog with M = N_R/N_d ratio typically 50x
%    - That's ~1 billion random points
%    - Need xi(r_p, pi) in 2D bins: 200 x 100 = 20,000 bins
%    - Need to run on ~1000-3000 mocks for covariance

% 2. COMPUTATIONAL BOTTLENECK IN EUCLID
%    - RR pair counting dominates when M >> 1 (M = 50)
%    - With split random trick: RR cost ~ M * N_d^2 / M^2 = N_d^2 / M
%    - DR cost ~ N_d * N_R = M * N_d^2
%    - DD cost ~ N_d^2
%    - Total dominated by DR for large M
%    - Uses kd-tree, octree, AND linked-list methods (3 independent implementations!)
%    - OpenMP parallelization (shared memory)

% 3. LINEAR CONSTRUCTION (LC) TRICK (Keihänen et al. 2022)
%    - Key insight: covariance of LS estimator can be decomposed into
%      terms that depend on M=1 and M=2 random catalogs only
%    - Construct covariance for arbitrary M from these two measurements
%    - Speedup factor of 14x for M=50 with 2 Mpc/h bins
%    - This is specifically for COVARIANCE estimation on mocks
%    - The actual xi measurement still needs M=50

% 4. WHAT kNN APPROACH CHANGES:
%    
%    a) RR is ANALYTIC (Erlang distribution) - zero cost
%       Euclid spends significant compute on RR even with split randoms
%       kNN: RR = 0 computation
%    
%    b) DR is O(N_R log N_d) not O(N_R * N_shell)
%       Each random does one kNN query, not a pair count across all bins
%       The kNN query traverses O(log N_d) nodes, not O(N_shell) pairs
%       For large r, N_shell ~ n_bar * 4pi r^2 dr can be very large
%    
%    c) DD is O(N_d log N_d) not O(N_d * N_shell) or O(N_d^2)
%       Same tree-traversal argument
%    
%    d) No binning choice for radial bins
%       Euclid uses 200 x 100 bins in (r_p, pi)
%       kNN gives continuous function, bin in post-processing if desired
%    
%    e) Covariance from dilution subsamples is FREE
%       No need for 1000-3000 mock runs
%       No need for LC trick
%       Built into the measurement

% 5. WHERE kNN DOES NOT HELP OR IS WORSE:
%    
%    a) The (r_p, pi) grid for large pi_max
%       Need large k_max or the lambda-metric trick
%       Standard pair counting in cylindrical bins may be simpler
%    
%    b) Very fine angular binning in mu
%       Each kNN neighbor contributes to one mu value
%       For 100 mu bins, you need enough neighbors per query point
%       The kNN approach has Poisson noise per (r, mu) cell
%    
%    c) Massive parallelization across nodes (MPI)
%       Euclid uses shared-memory OpenMP
%       kNN tree queries are embarrassingly parallel (each query independent)
%       so this is equally good for kNN
%    
%    d) The 2D (r_p, pi) measurement at fine binning
%       This is where conventional pair counting with cylindrical shells
%       may remain competitive, especially for Corrfunc-type optimized codes

% 6. REALISTIC COMPARISON FOR EUCLID NUMBERS:
%    N_d = 20 million (one redshift bin)
%    N_R = 50 * N_d = 1 billion
%    s_max = 200 Mpc/h, with 1 Mpc/h bins
%    
%    Standard LS (with tree acceleration):
%      DD: N_d * N_shell(s_max) ~ 2e7 * (n_bar * 4pi * 200^2 * 1) ~ 2e7 * 1e4 = 2e11
%      DR: N_R * N_shell ~ 1e9 * 1e4 = 1e13
%      RR: with split random, ~ N_R^2 / M = (1e9)^2 / 50 = 2e16 ... no, split reduces this
%      Actually RR with split: M subsets of N_d each, count within, ~ M * N_d^2 = 50 * (2e7)^2 = 2e16
%      This is huge. LC trick reduces covariance cost by 14x.
%    
%    kNN ladder:
%      Tree build: O(N_d log N_d) ~ 2e7 * 25 = 5e8 per dilution level
%      DR query: O(N_R log N_d) ~ 1e9 * 25 = 2.5e10 (for one lambda value)
%      DD query: O(N_d log N_d) ~ 5e8 per dilution level
%      RR: 0
%      
%      Total per level: ~2.5e10 (dominated by DR)
%      5 levels: ~1.3e11
%      
%      vs. standard DR alone: 1e13
%      Speedup: ~80x on DR term
%      
%      But: kNN query returns k=8 neighbors with full position info
%      Standard pair count bins into (r_p, pi) simultaneously
%      So per-query cost may differ by constant factors
%    
%    The real comparison needs wall-clock benchmarks, not op counts.
%    bosque at 10 ms for 1M queries on 100K data = 10 ns/query
%    Extrapolating: 1B queries on 20M data ~ 1B * (10 ns * log(20M)/log(100K)) ~ 15 seconds??
%    That can't be right. Need actual benchmarks.

% 7. WHAT TO WRITE IN THE PAPER:
%    - Replace the simplistic O(N^2) vs O(N log N) table
%    - Acknowledge the Euclid pipeline's actual optimizations:
%      split randoms, LC covariance, tree-accelerated pair counting
%    - Be honest about where kNN wins and where it doesn't
%    - The real win is: RR = 0, continuous in r, covariance from subsamples
%    - The marginal win is: DR cost reduction (log N vs N_shell, but N_shell
%      is also reduced by tree acceleration in standard codes)
%    - The new win is: free diagnostics (Delta_k, sigma_NL, DSC)
%    - Defer definitive benchmarks to the validation section / code paper
